%%%%%%%%%%%%%%%%%%%%%%%%%%%%%%%%%%%%%%%%%%%%%%%%%%%%%%%%%%%%%
%                                                           %
%                                                           %
%                                                           %
% COMIENZO DE ESCRITURA: 13 - MARZO - 2000                  %
% \'ULTIMA MODIFICACI\'ON: 5 - FEBRERO - 2002               %
%                                                           %
%                                                           %
%                                                           %
%%%%%%%%%%%%%%%%%%%%%%%%%%%%%%%%%%%%%%%%%%%%%%%%%%%%%%%%%%%%%
\documentclass{amsart}
%\renewcommand{\baselinestretch}{2} % DOBLE ESPACIO
%\renewcommand{\baselinestretch}{1.4}

\theoremstyle{plain}
\newtheorem*{theorem}{Theorem}
\newtheorem{lemma}{Lemma}

\newcommand{\supp}{\operatorname{supp}}
\newcommand{\spa}{\operatorname{span}}
\newcommand{\sgn}{\operatorname{sgn}}
\newcommand{\R}{\mathbb{R}}
\newcommand{\Z}{\mathbb{Z}}
\newcommand{\T}{\mathbb{T}}

%%%%%%%%%%%%%%%%%%%%%%%%%%%%%%%%%%%%%%%%%%%%%%%%%%%%%%%%%%%%%%%%%%%%%%%

\begin{document}

\title{Weak behaviour of Fourier-Neumann series}

\author[\'O. Ciaurri]{\'Oscar Ciaurri}
\address{Departamento de Matem\'aticas y Computaci\'on,
Universidad de La Rioja, Edificio J.~L.~Vives,
Calle Luis de Ulloa s/n, 26004 Logro\~no, Spain}
\email{oscar.ciaurri@dmc.unirioja.es}

\author[M. P\'erez]{Mario P\'erez}
\address{Departamento de Matem\'aticas,
Edificio de Matem\'aticas, Universidad de Zaragoza,
50009 Zaragoza, Spain}
\email{mperez@posta.unizar.es}

\author[J. L. Varona]{Juan~L.~Varona}
\address{Departamento de Matem\'aticas y Computaci\'on,
Universidad de La Rioja, Edificio J.~L.~Vives,
Calle Luis de Ulloa s/n, 26004 Logro\~no, Spain}
\email{jvarona@dmc.unirioja.es}
\urladdr{http://www.unirioja.es/dptos/dmc/jvarona/welcome.html}
\thanks{Research supported by grants of the DGI and UR}

\subjclass[2000]{Primary 42C10; Secondary 44A05}

\keywords{Bessel functions, Fourier series, Neumann series, Hankel transform.}

\thanks{PUBLISHED IN: Glasgow Math.\ J.\ \textbf{45} (2003), no.~1, 97--104;
and Glasgow Math.\ J.\ \textbf{45} (2003), no.~3, 567 (Corrigendum).}

\begin{abstract}
Let $J_\mu$ denote the Bessel function of order $\mu$. The functions
$x^{-\alpha-1}J_{\alpha+2n+1}(x)$, $n = 0, 1, 2, \dots$,
form an orthogonal system in the space $L^2((0,\infty),x^{2\alpha+1}dx)$
when $\alpha>-1$. In this paper we prove that the Fourier series associated
to this system is of restricted weak type for the endpoints of
the interval of mean convergence, while it is not of weak type if $\alpha
\geq 0$.
\end{abstract}

%%%%%%%%%%%%%%%%%%%%%%%%%%%%%%%%%%%%%%%%%%%%%%%%%%%%%%%%%%%%%%%%%%%%%%%
\maketitle
%%%%%%%%%%%%%%%%%%%%%%%%%%%%%%%%%%%%%%%%%%%%%%%%%%%%%%%%%%%%%%%%%%%%%%%

%%%%%%%%%%%%%%%%%%%%%%%%%%%%%%%%%%%%%%%%%%%%%%%%%%%%%%%%%%%%%%%%%%%%%%%
\section{Introduction and results}
%%%%%%%%%%%%%%%%%%%%%%%%%%%%%%%%%%%%%%%%%%%%%%%%%%%%%%%%%%%%%%%%%%%%%%%

Given a positive measure $\sigma$ on some space and an orthonormal system 
$\{\varphi_n\}_{n\geq 0}$ in $L^2(\sigma)$, the
Fourier series associated to $\{\varphi_n\}_{n\geq 0}$ is the sequence of
operators $S_n$ defined by
\[
   S_n f = \sum_{k=0}^n c_k(f) \varphi_k,
   \qquad f \in L^2(\sigma),
\]
where $c_k(f) = \int f \varphi_k \, d\sigma$.
The elementary property that
$\| S_n f - f \|_{L^2(\sigma)} \longrightarrow 0$
for every $f \in \overline{\spa{\{\varphi_n\}_{n\geq0}}}$
raises the same question with the $L^2(\sigma)$ norm replaced by the
$L^p(\sigma)$ norm, $1 \leq p \leq \infty$. By the Banach-Steinhaus theorem,
this is equivalent to the uniform boundedness
$\| S_n f \|_{L^p(\sigma)} \leq  C \| f \|_{L^p(\sigma)}$,
$f \in L^p(\sigma)$, $n \geq 0$.

Needless to say, the most important case is the trigonometric system on the
unit circle $\T$, for which the boundedness holds if $1 < p < \infty$
\cite{R}. For $p =\infty$ the answer is definitely negative, while for $p
=1$ the boundedness fails but there is a weak substitute in terms of the
Lorentz space $L^{1,\infty}(\T,d\theta)$:
\[
   \| S_n f \|_{L^{1,\infty}(\T,d\theta)} \leq C \| f \|_{L^1(\T,d\theta)},
   \qquad f \in L^1(\T,d\theta), \ n \geq 0.
\]
Here,
\begin{align*}
   \| f \|_{L^{p,\infty}(\sigma)}
   &= \sup_{y>0} y \lambda(y)^{1/p}
   = \| t^{1/p} f^*(t) \|_{L^\infty(\R^+,dt)},
   \qquad 1 \leq p < \infty
   \\
   \| f \|_{L^{p,r}(\sigma)}
   &= \left(\frac{r}{p} \int_0^\infty [t^{1/p} f^*(t)]^r \frac{dt}{t}
     \right)^{1/r}
   \qquad 1 \leq p < \infty, \ 1 \leq r < \infty,
\end{align*}
where $\lambda$ is the distribution function and $f^*$ the
nonincreasing rearrangement of $f$. There is a H\"older's inequality
$\| f \|_{p_1,p_2} \leq C \| f \|_{q_1,q_2} \| f \|_{r_1,r_2}$, $1/p_i =
1/q_i + 1/r_i$. Also, $\| f \|_{p,\infty} \leq C \| f \|_{p,p} = C \| f \|_p
\leq C_1 \| f \|_{p,1}$. The reader is referred to~\cite{H} or~\cite{SW} for
further details on $L^{p,r}$ spaces.

After the trigonometric system, the convergence of Fourier series has been
studied for a number of orthonormal systems, including Jacobi
polynomials~\cite{P1,P2,M1,Ch1,GPV1}, Hermite and Laguerre
polynomials~\cite{M2,M3}, generalized Jacobi polynomials~\cite{B} and Bessel
functions on $(0,1)$~\cite{GPRV}.

In~\cite{V,GPV2} the authors characterized the $L^p$ convergence for the
Fourier-Neumann series, i.e., the Fourier expansion associated to the
functions
\[
   \mathbf{j}_n^\alpha(x)
   = \sqrt{4n+2\alpha+2} \ J_{\alpha+2n+1}(x) x^{-\alpha-1},
   \qquad n = 0, 1, \dots
\]
which are orthonormal on $L^2((0,\infty), x^{2\alpha+1} \, dx)$
[$L^2(x^{2\alpha+1})$, from now on], see~\cite[\S~13.41~(7), p.~404]{W}
and~\cite[\S~13.42~(1), p.~405]{W}. Here, $\alpha > -1$ and $J_\nu$ is
the Bessel function of order $\nu$.

For each suitable function $f$, let $S_nf$ be the
$n$-th partial sum of its Fourier series with respect to the system
$\{\mathbf{j}_n^\alpha\}_{n=0}^\infty$, i.e.,
\begin{equation*}
   S_n(f,x) = \int_0^{\infty}f(t) K_n(x,t)t^{2\alpha+1}\,dt,
   \qquad
   K_n(x,t)=\sum_{k=0}^n \mathbf{j}_k^{\alpha}(x) \mathbf{j}_k^{\alpha}(t).
\end{equation*}
In this paper, we study the weak and restricted weak behaviour of these
series, i.e., the uniform boundedness
\[
   \| S_n f \|_{L^{p,\infty}(x^{2\alpha+1})}
   \leq  C \| f \|_{L^p(x^{2\alpha+1})},
   \qquad f \in L^p(x^{2\alpha+1}), \ n \geq 0
\]
or
\[
   \| S_n f \|_{L^{p,\infty}(x^{2\alpha+1})}
   \leq  C \| f \|_{L^{p,1}(x^{2\alpha+1})},
   \qquad f \in L^{p,1}(x^{2\alpha+1}), \ n \geq 0.
\]
Let us focus on the weak boundedness. The a priori assumption that
$\mathbf{j}_n^\alpha \in L^q(x^{2\alpha+1})$ ($n=0,1,\dots$, $1/p + 1/q =1$)
should be made so as to guarantee the existence of the Fourier coefficients
for any $f \in L^p(x^{2\alpha+1})$. Also, we must assume that
$\mathbf{j}_n^\alpha \in L^{p,\infty}(x^{2\alpha+1})$ if we want $S_n f$
to be in $L^{p,\infty}(x^{2\alpha+1})$. By Lemmas 1 and 2 below, these
assumptions hold if and only if $p_1 \leq p < p_2$, where $p_1 =
4(\alpha+1)/(2\alpha+3)$, $p_2 = 4(\alpha+1)/(2\alpha+1)$ if $\alpha \geq
0$, and $p_1=4/3$, $p_2 = 4$ if $-1 < \alpha < 0$. For the restricted weak
boundedness, the same arguments lead to the a priori assumptions that $p_1
\leq p \leq p_2$. Since the $L^p$-$L^p$ boundedness holds if and only if
$p_1 < p < p_2$, and it implies the $L^p$-$L^{p,\infty}$ and
$L^{p,1}$-$L^{p,\infty}$ boundedness, the real interest is in both endpoints.
We obtain a negative answer for the weak type in the case $\alpha \geq 0$,
and a positive answer for the restricted weak type.

\begin{theorem}
Let $\alpha \geq 0$, $p_1 = 4(\alpha+1)/(2\alpha+3)$, $p_2 =
4(\alpha+1)/(2\alpha+1)$. Then the partial sum operators $S_n$,
$n=0,1,\dots$, are not uniformly bounded as operators from
$L^{p_i}(x^{2\alpha+1})$ into $L^{p_i,\infty}(x^{2\alpha+1})$ but are
uniformly bounded as operators from $L^{p_i,1}(x^{2\alpha+1})$ into
$L^{p_i,\infty}(x^{2\alpha+1})$, $i=1,2$. In the case $-1 < \alpha < 0$ the
second statement holds with $p_1=4/3$ and $p_2=4$.
\end{theorem}

There is a close connection between $S_n$ and the Hankel
transform ${\mathcal H}_\alpha$ given by
\begin{equation}
\label{ec:hankel}
   {\mathcal H}_\alpha f(x)
   = \int_0^\infty \frac{J_\alpha(xy)}{(xy)^\alpha} f(y) y^{2\alpha+1} \, dy.
\end{equation}
It turns out from~\cite[\S~5.1~(8), p.~134]{W} and~\cite[proof of Lemma
4]{CGPV} that
$
   S_n f = {\mathcal H}_\alpha(\chi_{[0,1]}{\mathcal H}_\alpha f)
   - {\mathcal H}_{\alpha,n}(\chi_{[0,1]}{\mathcal H}_{\alpha,n} f)
   \equiv M_\alpha f - M_{\alpha,n}f,
$
where ${\mathcal H}_{\alpha,n}$ is given by~(\ref{ec:hankel}) with
$J_{\alpha+2n+2}$ in place of $J_\alpha$. Now, Kenig and Thomas~\cite{KT}
proved that the multiplier $M_\alpha$ is not bounded from
$L^p(x^{2\alpha+1})$ into $L^{p,\infty}(x^{2\alpha+1})$, $p =
4(\alpha+1)/(2\alpha+3)$. Chanillo~\cite{Ch2} proved that $M_\alpha$ is
bounded from $L^{p,1}(x^{2\alpha+1})$ into $L^{p,\infty}(x^{2\alpha+1})$, $p
= 4(\alpha+1)/(2\alpha+3)$. Some related uniform estimates were obtained by
Carbery, Romera and Soria~\cite{RS, CRS} in the context of the disc
multiplier.

Throughout this paper, unless otherwise stated, we use $C$, $C_1$
to denote positive constants independent of $n$ (and all other variables),
which can assume different values in different occurrences. As usual, we
write $f = O(g)$ in a given domain if $|f| \leq C g$. Finally, the standard
notation $a_+ = \max\{a,0\}$ will be used.

%%%%%%%%%%%%%%%%%%%%%%%%%%%%%%%%%%%%%%%%%%%%%%%%%%%%%%%%%%%%%%%%%%%%%%%%%%%
\section{Auxiliary results}
%%%%%%%%%%%%%%%%%%%%%%%%%%%%%%%%%%%%%%%%%%%%%%%%%%%%%%%%%%%%%%%%%%%%%

Some appropriate estimates for Bessel functions will be needed.
For instance,
\begin{align}
\label{ec:Jnu0}
   J_\nu(x) &= \frac{x^\nu}{2^\nu \Gamma(\nu+1)} + O(x^{\nu+2}),
   \quad  x\to 0+,
   \\
\label{ec:Jnuinfty}
   J_\nu(x) &= \sqrt{\frac2{\pi x}}
   \left[\cos\left(x-\frac{\nu\pi}2-\frac\pi4\right) + O(x^{-1})\right],
   \ \quad x\to \infty,
\end{align}
where the $O$ terms depend on $\nu$ (see~\cite[\S~3.1~(8),
p.~40]{W} and~\cite[\S~7.21~(1), p.~199]{W}). Some bounds for $J_\nu$ and
$J_\nu'$ with constants independent of $\nu$ are also available.
If $\nu > 0$, $0 < x \leq \nu/2$, and $a \geq -\nu$ then there exists some
constant $C_a$ depending only on $a$, such that
\begin{equation}
\label{ec:Jnuxa}
   |J_\nu(x)| x^a \le C_a \nu^{a-1/2} \left(\frac{e}{4}\right)^\nu
\end{equation}
(see~\cite[\S~3.31, p.~49]{W}). The formula $2 J_\nu' = J_{\nu-1} -
J_{\nu+1}$ proves the same bound for $J_\nu'(x)$, as well as the analogs
to~(\ref{ec:Jnu0}) and~(\ref{ec:Jnuinfty}):
\begin{align*}
   J_\nu'(x) &= \frac{x^{\nu-1}}{2^\nu \Gamma(\nu)} + O(x^{\nu+1}),
   \quad  x\to 0+,
   \\
   J_\nu'(x) &= \sqrt{\frac2{\pi x}}
   \left[-\sin\left(x-\frac{\nu\pi}2-\frac\pi4\right) + O(x^{-1})\right],
   \ \quad x\to \infty.
\end{align*}
It is easy to deduce from~(\ref{ec:Jnuxa}) and bounds done by Barcel\'o and
C\'ordoba (see~\cite[p.~661]{BC},~\cite[p.~24]{Co}) that
\begin{align}
\label{ec:Jnu}
   |J_\nu(x)| &\le C x^{-1/4} \left( |x - \nu| + \nu^{1/3} \right)^{-1/4},
   \qquad x \in (0,\infty),
   \quad \nu > 1
   \\
\label{ec:Jnu'}
   |J_\nu'(x)| &\le C x^{-3/4} \left( |x - \nu| + \nu^{1/3} \right)^{1/4},
   \qquad x \in (0,\infty),
   \quad \nu > 1
\end{align}
with some constant $C$ independent on $\nu$.
As a consequence, the following estimate for the norm of $x^b J_\nu(x)$ and
$x^b J_\nu'(x)$ in $L^p(x^{2\alpha+1})$ and $L^{p,\infty}(x^{2\alpha+1})$ can
be given.

\begin{lemma}
\label{lema:JnuLp}
Let $1 \leq p < \infty$, $\alpha > -1$, $b \in \mathbb{R}$ and $\nu > 1$. Let
$\lambda(4,\nu) = (\log \nu)^{1/4}$, $\lambda(p,\nu) = 1$ if $p \not= 4$.
Then
\begin{enumerate}
\item[(a)]
$x^b J_\nu(x) \in L^p(x^{2\alpha+1})$ if and only if
$p (b + \nu) + 2\alpha + 2 > 0$ and $p (b-\frac{1}{2}) + 2\alpha + 2 < 0$.
In that case,
$
   \| x^b J_\nu(x) \|_{L^p(x^{2\alpha+1})}
   \leq C \lambda(p,\nu)
   \nu^{2\frac{\alpha+1}{p} + b - \frac{1}{2}
   + \frac{2}{3} (\frac{1}{4} - \frac{1}{p})_+}.
$
\item[(b)]
$x^b J_\nu(x) \in L^{p,\infty}(x^{2\alpha+1})$ if and only if
$p (b + \nu) + 2\alpha + 2 \geq 0$ and $p (b-\frac{1}{2}) + 2\alpha + 2 \leq
0$.
In that case, 
$
   \| x^b J_\nu(x) \|_{L^{p,\infty}(x^{2\alpha+1})}
   \leq C \nu^{2\frac{\alpha+1}{p} + b - \frac{1}{2}
   + \frac{2}{3} (\frac{1}{4} - \frac{1}{p})_+}.
$
\item[(c)]
$x^b J_\nu'(x) \in L^p(x^{2\alpha+1})$ if and only if
$p (b + \nu - 1) + 2\alpha + 2 > 0$ and $p (b-\frac{1}{2}) + 2\alpha + 2 < 0$.
In that case, 
$
   \| x^b J_\nu'(x) \|_{L^p(x^{2\alpha+1})}
   \leq C \nu^{2\frac{\alpha+1}{p} + b - \frac{1}{2}}.
$
\item[(d)]
$x^b J_\nu'(x) \in L^{p,\infty}(x^{2\alpha+1})$ if and only if
$p (b + \nu - 1) + 2\alpha + 2 \geq 0$ and $p (b-\frac{1}{2}) + 2\alpha + 2
\leq 0$.
In that case,
$
   \| x^b J_\nu'(x) \|_{L^{p,\infty}(x^{2\alpha+1})}
   \leq C \nu^{2\frac{\alpha+1}{p} + b - \frac{1}{2}}.
$
\end{enumerate}
\end{lemma}

Similar results can be found in~\cite{BC, S}, so we will omit the proof
(details are given in~\cite[Chapter~2]{Ci}). Let us just mention that the
$L^p$ and $L^{p,\infty}$ conditions follow easily from~(\ref{ec:Jnu0}),
(\ref{ec:Jnuinfty}), and the analogs for $J_\nu'(x)$, while the norm estimates
are a consequence of~(\ref{ec:Jnuxa}),~(\ref{ec:Jnu}), and the analogs for
$J_\nu'(x)$.

Next lemma is the main step in the proof of the uniform restricted weak type:

\begin{lemma}
\label{lema:Lnu}
Let $\nu>1$, $1<p<\infty$, and
$L_\nu(f,x) = J_{\nu}(x^{1/2})H(t^{1/2} J'_{\nu}(t^{1/2})f(t),x)$,
where $H$ denotes the Hilbert transform.
Then, there exists a constant $C$, independent of $\nu$, such that
\begin{enumerate}
\item[(a)]
$\|L_\nu f\|_{L^p(dx)} \le C \|f\|_{L^p(dx)}$, $f\in L^p(dx)$, if $p<4$,
\item[(b)]
$\|L_\nu f\|_{L^{4,\infty}(dx)} \le C \|f\|_{L^{4,1}(dx)}$, $f\in
L^{4,1}(dx)$.
\end{enumerate}
\end{lemma}

\begin{proof}
(a) It follows from~(\ref{ec:Jnu}) that
\begin{equation*}
    \| L_\nu f \|_{L^p(dx)} \le C
    \| H(t^{1/2} J'_{\nu}(t^{1/2})f(t),x)
    \|_{L^p(x^{-p/8}(|x^{1/2}-\nu|+{\nu}^{1/3})^{-p/4})}.
\end{equation*}
Now, $x^{-p/8}(|x^{1/2}-\nu|+{\nu}^{1/3})^{-p/4} \in A_p$ uniformly in $\nu$
if $p<4$ (see~\cite{CGPV}, \cite{GPV2} or~\cite{V}). Thus, $H$ is a bounded
operator on $L^p(x^{-p/8}(|x^{1/2}-\nu|+{\nu}^{1/3})^{-p/4})$ and this,
together with~(\ref{ec:Jnu'}), proves (a).

(b) Let us write $L_\nu(f,x) = L_{\nu,1}(f,x) +
L_{\nu,2}(f,x)$, where
\begin{align*}
   L_{\nu,1}(f,x) 
   &= J_\nu(x^{1/2}) (|x^{1/2} - \nu| + \nu^{1/3})^{1/4}
   H \left(
   \frac{t^{1/2} J_\nu'(t^{1/2}) f(t)}{(|t^{1/2} - \nu| + \nu^{1/3})^{1/4}},
   x \right),
   \\
   L_{\nu,2}(f,x)
   &= J_\nu(x^\frac{1}{2}) H \left(
   \frac{t^\frac{1}{2} J_\nu'(t^\frac{1}{2}) f(t)
   [(|t^\frac{1}{2} - \nu| + \nu^\frac{1}{3})^\frac{1}{4} -
   (|x^\frac{1}{2} - \nu| + \nu^\frac{1}{3})^\frac{1}{4}]}
   {(|t^\frac{1}{2} -\nu| + \nu^\frac{1}{3})^\frac{1}{4}},
   x \right).
\end{align*}
The term $L_{\nu,1}(f,x)$ is easy to handle: we have
\[
   |J_\nu(x^{1/2})| (|x^{1/2} - \nu| + \nu^{1/3})^{1/4} \leq C x^{-1/8}
\]
and $x^{-1/2} \in A_4(0,\infty)$, so that
\begin{align*}
   \| L_{\nu,1}f \|_{L^{4,\infty}(dx)}
   &\leq C \left\| H \left(
   \frac{t^{1/2} J_\nu'(t^{1/2}) f(t)}{(|t^{1/2} - \nu| + \nu^{1/3})^{1/4}},
   x \right)
   \right\|_{L^4(x^{-1/2})}
   \\
   &\leq C \left\|
   \frac{x^{1/2} J_\nu'(x^{1/2}) f(x)}{(|x^{1/2} - \nu| + \nu^{1/3})^{1/4}}
   \right\|_{L^4(x^{-1/2})}
   \leq C \| f \|_{L^{4,1}(dx)}.
\end{align*}
Let us consider now $L_{\nu,2} f$. The elementary inequality $|a^{1/4} -
b^{1/4}| \leq a^{-3/4} |a - b|$, which holds for every $a, b > 0$, gives
\[
   \left|
   (|t^\frac{1}{2} - \nu| + \nu^\frac{1}{3})^\frac{1}{4} -
   (|x^\frac{1}{2} - \nu| + \nu^\frac{1}{3})^\frac{1}{4}
   \right|
   \leq t^{-1/2} (|t^\frac{1}{2} - \nu| + \nu^\frac{1}{3})^{-\frac{3}{4}}
   |t - x|.
\]
Thus,
\begin{align*}
   |L_{\nu,2}(f,x)|
   &\leq |J_\nu(x^{1/2})|
   \int_0^\infty |J_\nu'(t^{1/2})|
   (|t^\frac{1}{2} - \nu| + \nu^\frac{1}{3})^{-1} |f(t)| \, dt
   \\
   &\leq |J_\nu(x^{1/2})| \,
   \| J_\nu'(t^{1/2}) (|t^\frac{1}{2} - \nu| + \nu^\frac{1}{3})^{-1}
   \|_{L^{4/3,\infty}(dt)} \| f \|_{L^{4,1}(dt)}.
\end{align*}
Finally, Lemma~\ref{lema:JnuLp}(b) and a small change in
Lemma~\ref{lema:JnuLp}(d) give $\| J_\nu(x^\frac{1}{2})
\|_{L^{4,\infty}(dx)} \leq C$ and $\| J_\nu'(x^\frac{1}{2}) (|x^\frac{1}{2}
- \nu| + \nu^\frac{1}{3})^{-1} \|_{L^{4/3,\infty}(dx)} \leq C$.
\end{proof}

%%%%%%%%%%%%%%%%%%%%%%%%%%%%%%%%%%%%%%%%%%%%%%%%%%%%%%%%%%%%%%%%%%%%%%%
\section{Weak boundedness}
%%%%%%%%%%%%%%%%%%%%%%%%%%%%%%%%%%%%%%%%%%%%%%%%%%%%%%%%%%%%%%%%%%%%%%

Using that, for $\alpha >-1$,
\begin{multline*}
   \sum_{k=0}^n 2 (\alpha+2k+1) J_{\alpha+2k+1}(x) J_{\alpha+2k+1}(t)
   = \frac{xt}{x^2-t^2} \, [
   x J_{\alpha+1}(x) J_\alpha(t) - t J_\alpha(x) J_{\alpha+1}(t)
   \\
   + x J'_{\alpha+2n+2}(x) J_{\alpha+2n+2}(t)
   - t J_{\alpha+2n+2}(x) J'_{\alpha+2n+2}(t)
   ]
\end{multline*}
(see~\cite{GPV2, V}), we have
$S_n f = W_1 f - W_2 f + W_{3,n} f - W_{4,n} f$,
with
\begin{align*}
   W_1(f,x) &= \tfrac12 x^{-\alpha+1} J_{\alpha+1}(x)
   \, H\big( t^{\alpha/2} J_{\alpha}(t^{1/2}) f(t^{1/2}), x^2 \big),
   \\
   W_2(f,x) &= \tfrac12 x^{-\alpha} J_{\alpha}(x)
   \, H\big(t^{\alpha/2+1/2} J_{\alpha+1}(t^{1/2}) f(t^{1/2}), x^2 \big),
   \\
   W_{3,n}(f,x) &= \tfrac12 x^{-\alpha+1} J_\nu'(x)
   \, H\big(t^{\alpha/2} J_\nu(t^{1/2}) f(t^{1/2}), x^2 \big),
   \\
   W_{4,n}(f,x) &= \tfrac12 x^{-\alpha} J_\nu(x)
   \, H\big(t^{\alpha/2+1/2} J_\nu'(t^{1/2}) f(t^{1/2}), x^2 \big),
\end{align*}
and $\nu=\alpha+2n+2$. Here, $H$ denotes the Hilbert transform
on $(0,\infty)$.
The following was proved in~\cite[Theorem~1]{GPV2}:
\begin{gather}
\label{ec:acotW1}
   \| W_1 f \|_{L^p(x^{2\alpha+1})} \leq C \| f \|_{L^p(x^{2\alpha+1})},
   \quad
   \frac{2\alpha-1}{4(\alpha+1)} < \frac{1}{p} <
   \frac{2\alpha+3}{4(\alpha+1)};
   \\
\label{ec:acotW2}
   \| W_2 f \|_{L^p(x^{2\alpha+1})} \leq C \| f \|_{L^p(x^{2\alpha+1})},
   \quad
   \frac{2\alpha+1}{4(\alpha+1)} < \frac{1}{p} <
   \frac{2\alpha+5}{4(\alpha+1)};
   \\
\label{ec:acotW3n}
   \| W_{3,n} f \|_{L^p(x^{2\alpha+1})} \leq C \| f \|_{L^p(x^{2\alpha+1})},
   \quad
   \frac{2\alpha-\frac{1}{2}}{4(\alpha+1)} < \frac{1}{p} <
   \min\Big\{\frac{2\alpha+3}{4(\alpha+1)},\frac{3}{4}\Big\};
   \\
\label{ec:acotW4n}
   \| W_{4,n} f \|_{L^p(x^{2\alpha+1})} \leq C \| f \|_{L^p(x^{2\alpha+1})},
   \quad
   \max\Big\{\frac{2\alpha+1}{4(\alpha+1)},\frac{1}{4}\Big\}
   < \frac{1}{p} <
   \frac{2\alpha+\frac{9}{2}}{4(\alpha+1)}.
\end{gather}

Now, let $\alpha \geq 0$. As mentioned in the introduction, the $S_n$ are not
bounded operators from $L^p(x^{2\alpha+1})$ into $L^{p,\infty}(x^{2\alpha+1})$
if $p = 4(\alpha+1)/(2\alpha+1)$, so we only need to prove here that the
uniform weak boundedness fails for $p = 4(\alpha+1)/(2\alpha+3)$.

It follows from~(\ref{ec:acotW2}) and~(\ref{ec:acotW4n}) that $W_2$ and
$W_{4,n}$ are uniformly bounded from $L^p(x^{2\alpha+1})$ into itself, $p =
4(\alpha+1)/(2\alpha+3)$. Thus, it will be enough to find a sequence of
functions $\{f_n\}$ such that the inequality
\begin{equation}
\label{ec:W1+W3n}
   \| W_1 f_n + W_{3,n} f_n \|_{L^{p,\infty}(x^{2\alpha+1})}
   \leq C \| f_n \|_{L^p(x^{2\alpha+1})}
\end{equation}
fails for every constant $C$. Let
$
    f_n(t)
    = \sgn(J_{\alpha}(t))
    t^{-\frac{2\alpha+3}{2}}\chi_{[1,n]}(t)
$.
Then,
\begin{equation*}
   \| f_n \|_{L^p(x^{2\alpha+1})}
   = C (\log n)^\frac{1}{p}\, ,
   \qquad p = 4(\alpha+1)/(2\alpha+3).
\end{equation*}
Now, for $\nu = \alpha + 2n + 2$ and $x > 2 \nu$ we have 
\[
   |W_{3,n}(f_n,x)|
   \leq C x^{-\alpha-1} |J_\nu'(x)|
   \int_1^{n^2} t^{-\frac{3}{4}} |J_\nu(t^\frac{1}{2})| \, dt
   \leq C x^{-\alpha -\frac{3}{2}} \left(\frac{e}{4}\right)^{2n},
\]
where the last step follows from~(\ref{ec:Jnu'}) and~(\ref{ec:Jnuxa}). Thus,
\begin{equation}
\label{ec:W3nfn}
   \| \chi_{(2\nu,\infty)} W_{3,n} f_n
   \|_{L^{p,\infty}(x^{2\alpha+1})}
   \leq C \left(\frac{e}{4} \right)^{2n},
   \qquad p=4(\alpha+1)/(2\alpha+3).
\end{equation}
On the other hand, for $x > 2 \nu$ we have
\[
   |W_1(f_n,x)| 
   \geq C x^{-\alpha-1} |J_{\alpha+1}(x)|
   \int_1^{n^2} t^{-\frac{3}{4}}
   |J_\alpha(t^\frac{1}{2})| \, dt
   \geq C (\log n) x^{-\alpha-1}
   |J_{\alpha+1}(x)|,
\]
the last step following from~(\ref{ec:Jnuinfty}). Therefore,
\begin{equation}
\label{ec:W1fn}
   \| \chi_{(2\nu,\infty)}(x) W_1(f_n,x)
   \|_{L^{p,\infty}(x^{2\alpha+1})}
   \geq C \log n.
\end{equation}
Putting~(\ref{ec:W3nfn}) and~(\ref{ec:W1fn}) together, we get
\[
   \| W_1 f_n + W_{3,n} f_n
   \|_{L^{p,\infty}(x^{2\alpha+1})}
   \geq  C \log n,
   \qquad p = 4(\alpha+1)/(2\alpha+3)
\]
and~(\ref{ec:W1+W3n}) indeed fails. \qed

%%%%%%%%%%%%%%%%%%%%%%%%%%%%%%%%%%%%%%%%%%%%%%%%%%%%%%%%%%%%%%%%%%%
\section{Restricted weak boundedness}
%%%%%%%%%%%%%%%%%%%%%%%%%%%%%%%%%%%%%%%%%%%%%%%%%%%%%%%%%%%%%%%%%%%

By duality, we only need to prove that the self adjoint operators $S_n$ are
uniformly of restricted weak type in two cases: $\alpha \geq 0$, $p =
\frac{4(\alpha+1)}{2\alpha+1}$, and $-1 < \alpha < 0$, $p = 4$.

{\it Case $\alpha \ge0$ and $p = \frac{4(\alpha+1)}{2\alpha+1}$.}
From~(\ref{ec:acotW1}) and~(\ref{ec:acotW3n}), we conclude that $W_1$ and
$W_{3,n}$ are uniformly bounded from $L^p(x^{2\alpha+1})$ into itself.
Therefore, it is enough to prove that $W_2$ and $W_{4,n}$ are uniformly
bounded from
$L^{p,1}(x^{2\alpha+1})$ into $L^{p,\infty}(x^{2\alpha+1})$, i.e.,
\[
   \| W_{4,n} f \|_{L^{p,\infty}(x^{2\alpha+1})}
   \leq C \| f \|_{L^{p,1}(x^{2\alpha+1})},
   \ f \in L^{p,1}(x^{2\alpha+1})
\]
and the same inequality for $W_2$. We will consider only $W_{4,n}$,
since the boundedness of $W_2$ is completely analogous. Let $f \in
L^{p,1}(x^{2\alpha+1})$ and, for each $k \in \Z$,
\[
   I_k = [2^k, 2^{k+1}),
   \qquad
   f_1^k = f \chi_{(0,2^{k-1}) \cup [2^{k+2},\infty)},
   \qquad
   f_2^k = f \chi_{[2^{k-1},2^{k+2})}.
\]
Thus, $f = f_1^k + f_2^k$ for each $k \in \Z$ and
\begin{equation}
\label{ec:descW4n}
   |W_{4,n} (f,x)| 
   \leq \sum_{k \in \Z} |W_{4,n} (f_1^k,x)| \chi_{I_k}(x)
   + \sum_{k \in \Z} |W_{4,n} (f_2^k,x)| \chi_{I_k}(x).
\end{equation}
Let $x \in I_k$. Then, it is easy to
check that $|x^2 - y^2| \geq \frac{3}{4} y^2$ if $y \in (0,2^{k-1}) \cup
[2^{k+2},\infty)$. Hence, after a change of variable we get
\[
   |W_{4,n} (f_1^k,x)| 
   \leq C x^{-\alpha} |J_\nu(x)|
   \int_0^\infty y^{\alpha}
   |J_\nu'(y)| |f(y)| \, dy
\]
and
\[
   \sum_{k \in \Z} |W_{4,n} (f_1^k,x)| \chi_{I_k}(x)
   \leq C x^{-\alpha} |J_\nu(x)|\,
   \| x^{-\alpha-1} J_\nu'(x)
   \|_{L^{q,\infty}(x^{2\alpha+1})}
   \| f \|_{L^{p,1}(x^{2\alpha+1})},
\]
where $\frac{1}{p} + \frac{1}{q} = 1$. Therefore, the first term
in~(\ref{ec:descW4n}) is bounded:
\[
   \| \sum_{k \in \Z} |W_{4,n} (f_1^k,x)| \chi_{I_k}(x)
   \|_{L^{p,\infty}(x^{2\alpha+1})}
   \leq C \| f \|_{L^{p,1}(x^{2\alpha+1})},
\]
by Lemma~\ref{lema:JnuLp}(b) and (d).
Let us consider
now the second term. If $x \in I_k$,
\[
   |W_{4,n} (f_2^k,x)| 
   \leq C
   2^{-\alpha k} \left|L_\nu(t^\frac{\alpha}{2} f_k^2(t^{1/2}), x^2)
   \right|.
\]
From Lemma~\ref{lema:Lnu} it follows that
\begin{align*}
   \|W_{4,n} (f_2^k,x) \chi_{I_k}(x)\|_{L^{p,\infty}(x^{2\alpha+1})}    
   \leq C
   2^{-\alpha k + \frac{2\alpha k}{p}}
   \|  \chi_{I_k}(x) L_\nu(t^\frac{\alpha}{2} f_k^2(t^{1/2}), x^2)
   \|_{L^{p,\infty}(x\, dx)}
   \\
   \leq C
   2^{-\alpha k + \frac{2\alpha k}{p}}
   \| x^\frac{\alpha}{2} f_k^2(x^{1/2}) \|_{L^{p,1}(dx)}
   \leq C \| f \chi_{[2^{k-2},2^{k+2})} \|_{L^{p,1}(x^{2\alpha+1})}.
\end{align*}
Then,
\begin{equation}
\label{ec:W4nf2k}
   \| \sum_{k \in \Z} |W_{4,n} (f_2^k,x)|
   \chi_{I_k}(x)\|_{L^{p,\infty}(x^{2\alpha+1})} 
   \leq C \| f \|_{L^{p,1}(x^{2\alpha+1})}.
\end{equation}

{\it Case $-1 < \alpha < 0$ and $p =4$.} Now, $W_1$, $W_2$, and $W_{3,n}$
are uniformly bounded from $L^4(x^{2\alpha+1})$ into itself
(see~(\ref{ec:acotW1}), (\ref{ec:acotW2}), (\ref{ec:acotW3n})), so we only
need to prove that
\[
   \| W_{4,n} f \|_{L^{4,\infty}(x^{2\alpha+1})}
   \leq C \| f \|_{L^{4,1}(x^{2\alpha+1})},
   \ f \in L^{4,1}(x^{2\alpha+1}).
\]
The above proof of~(\ref{ec:W4nf2k}) remains valid, while only minor
changes are necessary for the first term in~(\ref{ec:descW4n}): it is not
difficult to check that
\[
   \frac{y^{2+\frac{\alpha}{2}} x^{-\frac{\alpha}{2}}}{|x^2-y^2|}
   \leq \frac{4}{3}
\]
if $x \in I_k$ and $y \in (0,2^{k-1}) \cup [2^{k+2},\infty)$. Then, it
follows that
\[
   |W_{4,n} (f_1^k,x)| \leq C x^{-\frac{\alpha}{2}} |J_\nu(x)|
   \int_0^\infty y^{\frac{\alpha}{2}} |J_\nu'(y)| \, |f(y)|
   \,dy.
   \qed
\]

%%%%%%%%%%%%%%%%%%%%%%%%%%%%%%%%%%%%%%%%%%%%%%%%%%%%%%%%%%%%%%%%%%%%%%%%%

%\bibliographystyle{amsplain}

\begin{thebibliography}{99}

\bibitem{B}
V. M. Badkov,
{Convergence in the mean and almost everywhere of Fourier series in
polynomials orthogonal on an interval},
\textit{Math. USSR Sb.}
\textbf{24} (1974), 223--256.

\bibitem{BC} J. A. Barcel\'o and A. C\'ordoba,
{Band-limited functions: $L^p$-convergence},
\textit{Trans. Amer. Math. Soc.}
\textbf{313} (1989), 655--669.

\bibitem{CRS} A. Carbery, E. Romera and F. Soria,
{Radial weights and mixed norm inequalities for the disc multiplier},
\textit{J. Funct. Anal.}
\textbf{109} (1992), 52--75.

\bibitem{Ch1} S. Chanillo,
{On the weak behaviour of partial sums of Legendre series},
\textit{Trans. Amer. Math. Soc.}
\textbf{268} (1981), 367--376.

\bibitem{Ch2} S. Chanillo,
{The multiplier for the ball and radial functions},
\textit{J. Funct. Anal.}
\textbf{55} (1984), 18--24.

\bibitem{Ci} \'O. Ciaurri,
{\it Aproximaci\'on de funciones cuya transformada de Hankel est\'a soportada
en $[0,1]$,}
Ph.~D.~Thesis, University of La Rioja, Logro\~no (Spain), 2000.

\bibitem{CGPV} \'O. Ciaurri, J. J. Guadalupe, M. P\'erez and J. L. Varona,
{Mean and almost everywhere convergence of Fourier-Neumann series},
\textit{J. Math. Anal. Appl.}
\textbf{236} (1999), 125--147.

\bibitem{Co} A. C\'ordoba,
{The disc multiplier},
\textit{Duke Math. J.}
\textbf{58} (1989), 21--29.

%\bibitem{EMOT} A. Erd\'elyi, W. Magnus, F. Oberhettinger and F. Tricomi,
%{\it Higher Transcendental Functions,} Vol.~II,
%McGraw-Hill, New York, 1953.

\bibitem{GPRV} J. J. Guadalupe, M. P\'erez, F. J. Ruiz and J. L. Varona,
{Mean and weak convergence of Fourier-Bessel series},
\textit{J. Math. Anal. Appl.}
\textbf{173} (1993), 370--389.

\bibitem{GPV1} J. J. Guadalupe, M. P\'erez and J. L. Varona,
{Weak behaviour of Fourier-Jacobi series},
\textit{J. Approx. Theory}
\textbf{61} (1990), 222--238.

\bibitem{GPV2} J. J. Guadalupe, M. P\'erez and J. L. Varona,
{Commutators and analytic dependence of Fourier-Bessel series on
$(0,\infty)$},
\textit{Canad. Math. Bull.}
\textbf{42} (1999), 198--208.

\bibitem{H} G. H. Hunt,
{On $L(p,q)$ spaces},
\textit{Enseign. Math.}
\textbf{12} (1966), 249--276.

\bibitem{KT} C. E. Kenig and P. A. Tomas,
{The weak behaviour of spherical means},
\textit{Proc. Amer. Math. Soc.}
\textbf{78} (1980), 48--50.

\bibitem{M1} B. Muckenhoupt,
{Mean convergence of Jacobi series},
\textit{Proc. Amer. Math. Soc.}
\textbf{23} (1969), 306--310.

\bibitem{M2} B. Muckenhoupt,
{Mean convergence of Hermite and Laguerre series.~I},
\textit{Trans. Amer. Math. Soc.}
\textbf{147} (1970), 419--431.

\bibitem{M3} B. Muckenhoupt,
{Mean convergence of Hermite and Laguerre series.~II},
\textit{Trans. Amer. Math. Soc.}
\textbf{147} (1970), 433--460.

\bibitem{P1} H. Pollard,
{The mean convergence of orthogonal series.~II},
\textit{Trans. Amer. Math. Soc.}
\textbf{63} (1948), 355--367.

\bibitem{P2} H. Pollard,
{The mean convergence of orthogonal series.~III},
\textit{Duke Math. J.}
\textbf{16} (1949), 189--191.

\bibitem{R} M. Riesz,
{Sur les fonctions conjugu\'ees},
\textit{Math. Z.}
\textbf{27} (1927), 218--244.

\bibitem{RS} E. Romera and F. Soria,
{Endpoint estimates for the maximal operator associated to
spherical partial sums on radial functions},
\textit{Proc. Amer. Math. Soc.}
\textbf{111} (1991), 1015--1022.

\bibitem{SW} E. M. Stein and G. Weiss,
{\it Introduction to Fourier Analysis on Euclidean Spaces,}
Princeton Univ. Press, Princeton, NJ, 1975.

\bibitem{S} K. Stempak,
{A weighted uniform $L^p$-estimate of Bessel functions:
A note on a paper of Guo},
\textit{Proc. Amer. Math. Soc.}
\textbf{128} (2000), 2943--2945.

\bibitem{V} J. L. Varona,
{Fourier series of functions whose Hankel transform is supported on $[0,1]$},
\textit{Constr. Approx.}
\textbf{10} (1994), 65--75.

\bibitem{W} G. N. Watson,
{\it A Treatise on the Theory of Bessel Functions} (2nd edition),
Cambridge Univ. Press,
Cambridge, 1944.

\end{thebibliography}


\end{document}